\section{Appendix}

\subsection{Prompt Used for Final ECG Question Answering}
\label{app:final_prompt}

The final adapter-based ECG question-answering system used a fixed textual prompt template. For each ECG-QA example, the question text was inserted into the following prompt:

\begin{verbatim}
You are a cardiology assistant.
Answer the ECG question using the ECG representation.
Question: {question}
Answer:
\end{verbatim}

The ECG recording itself was not inserted into the prompt as natural-language text. Instead, the precomputed CSFM ECG embedding was passed through the trained adapter to produce eight continuous ECG-derived soft prompt embeddings. These soft prompt embeddings were prepended to the token embeddings of the textual prompt before being supplied to the frozen language model.

During forced-choice evaluation, the model separately scored the two candidate completions ``yes'' and ``no''. The predicted answer was the candidate assigned the higher likelihood by the model.

\subsection{Additional Adapter Model Comparison Results}
\label{app:model_comparison}

Table~\ref{tab:appendix_model_comparison} reports additional validation results for adapter variants and training strategies that were used during model development.

\begin{table}[H]
\centering
\caption{Additional adapter model comparison results on the validation set.}
\label{tab:appendix_model_comparison}
\small
\begin{tabular}{p{0.42\linewidth}cccc}
\hline
\textbf{Experiment} & \textbf{Accuracy} & \textbf{Bal. acc.} & \textbf{Macro-F1} & \textbf{Yes / no recall} \\
\hline
3B adapter, original training set & 0.790 & 0.738 & 0.736 & 0.626 / 0.850 \\
3B adapter, class-weighted loss & 0.788 & 0.735 & 0.733 & 0.620 / 0.851 \\
3B adapter + LoRA, initial run & 0.803 & 0.730 & 0.739 & 0.570 / 0.890 \\
3B adapter + LoRA, resumed run & 0.791 & 0.738 & 0.736 & 0.624 / 0.852 \\
3B adapter, 40\% yes upsampling & 0.765 & 0.757 & 0.729 & 0.738 / 0.775 \\
8B adapter, original training set & 0.794 & 0.744 & 0.741 & 0.634 / 0.853 \\
8B adapter, 40\% yes upsampling & 0.789 & 0.764 & 0.748 & 0.710 / 0.819 \\
8B adapter, 50\% yes upsampling & 0.756 & 0.764 & 0.726 & 0.781 / 0.747 \\
\hline
\end{tabular}
\end{table}
\normalsize

\subsection{Clinical Group Results}
\label{app:clinical_group_results}

Table~\ref{tab:appendix_clinical_group_results} gives the exact validation values for the clinical-group comparison between the Llama-3.1-8B-Instruct adapter trained on the original training set and the same model trained with 40\% positive-example upsampling.

\begin{table}[H]
\centering
\caption{Clinical-group validation results before and after 40\% positive-example upsampling.}
\label{tab:appendix_clinical_group_results}
\small
\begin{tabular}{lrrrrr}
\hline
\textbf{Clinical group} & \textbf{N} & \textbf{Yes} & \textbf{No} & \textbf{Original bal. acc.} & \textbf{Upsampled bal. acc.} \\
\hline
Rhythm & 399 & 119 & 280 & 0.857 & 0.872 \\
Conduction & 387 & 117 & 270 & 0.782 & 0.840 \\
Normal & 60 & 20 & 40 & 0.800 & 0.762 \\
Hypertrophy & 300 & 90 & 210 & 0.786 & 0.729 \\
Form/morphology & 2411 & 651 & 1760 & 0.715 & 0.751 \\
ST/T change & 1073 & 258 & 815 & 0.711 & 0.726 \\
MI/injury & 159 & 39 & 120 & 0.860 & 0.727 \\
\hline
\end{tabular}
\end{table}
\normalsize

\subsection{Full Per-Label Results}
\label{app:full_per_scp_results}

Table~\ref{tab:appendix_full_per_scp_results} reports the full selected-label validation comparison used to support the per-SCP analysis. These labels correspond to SCP-code attributes from the filtered ECG-QA task, but are shown using clinical descriptions for readability.

% This table requires \usepackage{longtable} in the preamble.
\small
\begin{longtable}{p{0.36\linewidth}lrrr}
\caption{Full per-label validation results before and after 40\% positive-example upsampling.}\label{tab:appendix_full_per_scp_results}\\
\hline
\textbf{Clinical label} & \textbf{Group} & \textbf{N} & \textbf{Original bal. acc.} & \textbf{Upsampled bal. acc.} \\
\hline
\endfirsthead
\hline
\textbf{Clinical label} & \textbf{Group} & \textbf{N} & \textbf{Original bal. acc.} & \textbf{Upsampled bal. acc.} \\
\hline
\endhead
Complete left bundle branch block & Conduction & 51 & 0.963 & 0.975 \\
Complete right bundle branch block & Conduction & 60 & 0.963 & 0.963 \\
First degree AV block & Conduction & 60 & 0.713 & 0.738 \\
Incomplete left bundle branch block & Conduction & 18 & 0.800 & 0.933 \\
Incomplete right bundle branch block & Conduction & 60 & 0.762 & 0.788 \\
Left anterior fascicular block & Conduction & 60 & 0.762 & 0.887 \\
Left posterior fascicular block & Conduction & 18 & 0.500 & 0.800 \\
Non-specific intraventricular conduction disturbance (block) & Conduction & 60 & 0.625 & 0.688 \\
Abnormal QRS & Form/morphology & 60 & 0.675 & 0.662 \\
Atrial premature complex & Form/morphology & 60 & 0.812 & 0.775 \\
High QRS voltage & Form/morphology & 60 & 0.660 & 0.670 \\
Inverted T-waves & Form/morphology & 501 & 0.699 & 0.792 \\
Low QRS voltages in the frontal and horizontal leads & Form/morphology & 107 & 0.702 & 0.746 \\
Non-specific ST depression & Form/morphology & 622 & 0.715 & 0.708 \\
Non-specific ST elevation & Form/morphology & 30 & 0.420 & 0.500 \\
Non-specific T-wave changes & Form/morphology & 364 & 0.732 & 0.795 \\
Prolonged PR interval & Form/morphology & 60 & 0.762 & 0.812 \\
Q waves present & Form/morphology & 421 & 0.688 & 0.735 \\
T-wave abnormality & Form/morphology & 66 & 0.464 & 0.500 \\
Ventricular premature complex & Form/morphology & 60 & 0.900 & 0.938 \\
Left atrial overload/enlargement & Hypertrophy & 52 & 0.858 & 0.675 \\
Left ventricular hypertrophy & Hypertrophy & 236 & 0.779 & 0.730 \\
Right atrial overload/enlargement & Hypertrophy & 12 & 0.700 & 0.900 \\
Anteroseptal myocardial infarction & MI/injury & 60 & 0.850 & 0.750 \\
Subendocardial injury in anterolateral leads & MI/injury & 51 & 0.872 & 0.669 \\
Subendocardial injury in anteroseptal leads & MI/injury & 42 & 0.900 & 0.757 \\
Subendocardial injury in lateral leads & MI/injury & 6 & 0.500 & 0.500 \\
Normal ECG & Normal & 60 & 0.800 & 0.762 \\
Atrial fibrillation & Rhythm & 60 & 0.925 & 0.962 \\
Atrial flutter & Rhythm & 36 & 0.983 & 0.967 \\
Normal functioning artificial pacemaker & Rhythm & 51 & 1.000 & 1.000 \\
Sinus arrhythmia & Rhythm & 60 & 0.787 & 0.725 \\
Sinus bradycardia & Rhythm & 60 & 0.863 & 0.863 \\
Sinus rhythm & Rhythm & 60 & 0.600 & 0.725 \\
Sinus tachycardia & Rhythm & 60 & 0.963 & 0.900 \\
Supraventricular tachycardia & Rhythm & 12 & 0.750 & 1.000 \\
Ischemic in anterior leads & ST/T change & 18 & 0.467 & 0.500 \\
Ischemic in anterolateral leads & ST/T change & 60 & 0.675 & 0.713 \\
Ischemic in anteroseptal leads & ST/T change & 36 & 0.900 & 0.817 \\
Ischemic in inferior leads & ST/T change & 30 & 0.480 & 0.580 \\
Ischemic in inferolateral leads & ST/T change & 36 & 0.833 & 0.850 \\
Ischemic in lateral leads & ST/T change & 48 & 0.537 & 0.500 \\
Long QT-interval & ST/T change & 30 & 0.680 & 0.720 \\
Non-diagnostic T abnormalities & ST/T change & 455 & 0.768 & 0.762 \\
Non-specific ischemic & ST/T change & 60 & 0.850 & 0.938 \\
Non-specific ST changes & ST/T change & 300 & 0.581 & 0.610 \\
\hline
\end{longtable}
\normalsize

\subsection{Abbreviations}
\label{app:abbreviations}

\begin{table}[H]
\centering
\caption{Abbreviations used in the dissertation.}
\label{tab:appendix_abbreviations}
\begin{tabular}{ll}
\hline
\textbf{Abbreviation} & \textbf{Meaning} \\
\hline
ECG & Electrocardiogram \\
LLM & Large language model \\
CSFM & Cardiac Sensing Foundation Model \\
ECG-QA & Electrocardiogram question answering \\
SCP & Standard communications protocol for computer-assisted electrocardiography \\
AUC & Area under the curve \\
ROC-AUC & Receiver operating characteristic area under the curve \\
F1 & Harmonic mean of precision and recall \\
LoRA & Low-rank adaptation \\
SLURM & Simple Linux Utility for Resource Management \\
\hline
\end{tabular}
\end{table}
